\section{Objetivos específicos}

\begin{enumerate}
    \item Identificar, mediante el uso de modelos estadísticos y técnicas de ciencia de datos aplicadas a información transaccional parcial, las variables transaccionales y de perfilamiento del cliente que resultan relevantes para explicar la \emph{participación en la billetera} (SoW)) y el \emph{tamaño de la billetera} (SioW) de los clientes empresariales de la empresa focal, considerando explícitamente las restricciones derivadas de la ausencia de información externa sobre competidores.

    \item Comparar, dentro de un marco de modelamiento estadístico basado en ciencia de datos, distintas especificaciones del conjunto de covariables transaccionales empleadas en el modelo \emph{pogit}, incluyendo una especificación basada en variables originales y una alternativa basada en reducción de dimensionalidad mediante Análisis de Componentes Principales (PCA), con el fin de evaluar su impacto sobre la estabilidad, parsimonia e interpretabilidad de las estimaciones de SoW y SioW.

    \item Analizar la estructura de segmentación de los clientes empresariales obtenida mediante técnicas de aprendizaje no supervisado (k-means y k-medoids) y su correspondencia exploratoria con los patrones emergentes en el espacio conjunto SoW--SioW estimado por el modelo \emph{pogit}, como apoyo a la interpretación económica de los resultados y a la toma de decisiones estratégicas de marketing en la empresa focal.
\end{enumerate}


